\documentclass[12pt,a4paper]{article}
\usepackage[utf8]{inputenc}
\usepackage[french]{babel}
\usepackage[T1]{fontenc}
\usepackage{amsmath,amssymb}
\usepackage{graphicx}
\usepackage{tcolorbox}
\usepackage{fancyhdr}
\usepackage{enumitem}
\usepackage{geometry}
\usepackage{listings}
\usepackage{tikz}
\usepackage{hyperref}

\geometry{top=2.5cm,bottom=2.5cm,left=2.5cm,right=2.5cm}

\hypersetup{
    colorlinks=true,
    linkcolor=blue,
    filecolor=magenta,      
    urlcolor=cyan,
    pdftitle={Cours Complet: Cohérence des Caches & MEF Snoopy MSI},
}

\tcbuselibrary{skins,breakable}

% Styles de boîtes pédagogiques
\newtcolorbox{definitionbox}[1]{
  colback=blue!5!white,
  colframe=blue!75!black,
  fonttitle=\bfseries,
  title=📖 Définition Simple : #1,
  breakable
}

\newtcolorbox{analogybox}[1]{
  colback=amber!10!white,
  colframe=amber!85!black,
  fonttitle=\bfseries,
  title=💡 Analogie de la Vie Quotidienne : #1,
  breakable
}

\newtcolorbox{examplebox}[1]{
  colback=green!5!white,
  colframe=green!60!black,
  fonttitle=\bfseries,
  title=💻 Exemple Concret : #1,
  breakable
}

\newtcolorbox{transitionbox}[1]{
  colback=purple!5!white,
  colframe=purple!75!black,
  fonttitle=\bfseries,
  title=🔄 Transition MEF : #1,
  breakable
}

\newtcolorbox{alertbox}[1]{
  colback=red!5!white,
  colframe=red!75!black,
  fonttitle=\bfseries,
  title=⚠️ Piège de Cours / Erreur Fréquente : #1,
  breakable
}

\pagestyle{fancy}
\fancyhf{}
\rhead{Systèmes Multiprocesseurs \& Machines Parallèles}
\lhead{MEF Snoopy Cache (Protocole MSI)}
\cfoot{\thepage}

\title{\textbf{\Huge Cours Magistral d'Introduction Zéro-Prérequis}\\[0.5em]
\Large Cohérence des Mémoires Caches \& Machine à États Finis (MEF) Snoopy MSI}
\author{\textbf{Professeur Virtuel Antigravity} \\ Cours de Systèmes à Multiprocesseurs et Machines Parallèles}
\date{\today}

\begin{document}

\maketitle
\tableofcontents
\newpage

% ==========================================
\section{Étape 1 — Identifier le Domaine}
% ==========================================

\subsection{De quel domaine informatique parle cette image ?}
Cette étude relève de l'**Architecture des Ordinateurs** et plus particulièrement des **Systèmes Multiprocesseurs et Machines Parallèles**. C'est la discipline informatique qui étudie comment concevoir le matériel (ordinateurs, processeurs, puces électroniques) pour faire fonctionner plusieurs moteurs de calcul en même temps afin d'exécuter des programmes de manière ultra-rapide.

\subsection{À quoi sert ce domaine ?}
Ce domaine sert à augmenter massivement la puissance de calcul des ordinateurs modernes. Au lieu d'attendre qu'un seul processeur devienne deux fois plus rapide (ce qui est technologiquement devenu très difficile à cause des limites physiques et de la chaleur dégagée), les ingénieurs mettent **plusieurs processeurs (ou cœurs)** dans une même machine pour accomplir plusieurs tâches en parallèle.

\subsection{Pourquoi existe-t-il ?}
Historiquement, la vitesse des processeurs augmentait chaque année (Loi de Moore). Mais vers 2004, nous avons atteint un mur physique : augmenter la fréquence (la vitesse de l'horloge interne) chauffait tellement la puce qu'elle risquait de fondre. La solution a été d'adopter le **calcul parallèle** : au lieu d'un seul ouvrier très rapide, nous employons une équipe d'ouvriers travaillant en même temps.

\subsection{Quels problèmes cherche-t-il à résoudre ?}
Le grand défi du calcul parallèle est la **coordination** et le **partage des informations**. Lorsque plusieurs ouvriers (processeurs) travaillent sur un même projet et consultent ou modifient les mêmes documents (la mémoire), il faut garantir que :
\begin{enumerate}
    \item Tout le monde voit les informations à jour.
    \item Deux ouvriers ne modifient pas un document dans leur coin sans que les autres ne soient avertis.
\end{enumerate}
C'est précisément le problème de la **Cohérence de Cache** que résout le schéma étudié.

\newpage
% ==========================================
\section{Étape 2 — Tous les Prérequis (Décomposition Notions)}
% ==========================================

Avant d'aborder le schéma, décortiquons **chaque terme technique** à partir de zéro.

% 1. Processeur (CPU)
\subsection{Processeur (CPU - Central Processing Unit)}
\begin{definitionbox}{Processeur (CPU)}
Le cerveau ou le moteur de calcul de l'ordinateur. C'est lui qui exécute les instructions des programmes (calculs mathématiques, comparaisons, décisions).
\end{definitionbox}

\begin{analogybox}{Le Processeur}
Considérez le processeur comme le **chef cuisinier** dans un restaurant. C'est lui qui lit les recettes et prépare activement les plats.
\end{analogybox}

\begin{examplebox}{Le Processeur}
Quand vous cliquez sur un bouton dans un jeu vidéo, le CPU calcule la nouvelle position de votre personnage et met à jour son score.
\end{examplebox}
\textbf{Pourquoi il existe :} Sans CPU, l'ordinateur serait une machine morte sans capacité d'exécution.\\
\textbf{Comment il est utilisé :} Il lit les instructions depuis la mémoire, exécute les opérations puis sauvegarde les résultats.

% 2. Mémoire / RAM
\subsection{Mémoire Principale (RAM - Random Access Memory)}
\begin{definitionbox}{Mémoire Principale (RAM)}
L'espace de stockage de travail temporaire, rapide mais grand, où sont conservées les données et programmes en cours d'exécution.
\end{definitionbox}

\begin{analogybox}{La Mémoire Principale}
C'est la **grande étagère du garde-manger** située au fond de la cuisine. Elle contient tous les ingrédients (données), mais le cuisinier doit se déplacer pour y aller.
\end{analogybox}

\begin{examplebox}{La Mémoire Principale}
Une barrette de RAM DDR4 ou DDR5 de 16 Go intégrée dans votre ordinateur portable.
\end{examplebox}
\textbf{Pourquoi elle existe :} Le disque dur (SSD) est trop lent pour fournir les données directement au CPU pendant un calcul. La RAM est beaucoup plus rapide.\\
\textbf{Comment elle est utilisée :} Elle stocke les variables de vos programmes ouvertes à un instant $t$.

% 3. Mémoire Cache
\subsection{Mémoire Cache}
\begin{definitionbox}{Mémoire Cache}
Une toute petite mémoire ultra-rapide située directement à côté (ou à l'intérieur) du CPU pour garder en main les données fréquemment utilisées.
\end{definitionbox}

\begin{analogybox}{La Mémoire Cache}
C'est le **plan de travail du cuisinier**. Au lieu d'aller chercher le sel au garde-manger (RAM) à chaque pincée, il pose le pot de sel directement sur son plan de travail (Cache).
\end{analogybox}

\begin{examplebox}{La Mémoire Cache}
Le cache L1 de 32 Ko intégré directement dans un cœur de processeur Intel Core i7.
\end{examplebox}
\textbf{Pourquoi elle existe :} La RAM est devenue trop lente par rapport à l'extrême vitesse du CPU (problème appelé \textit{Memory Wall}). Sans cache, le CPU passerait 90\% de son temps à attendre la RAM.\\
\textbf{Comment elle est utilisée :} Lorsque le CPU demande une donnée, il regarde d'abord dans son cache. Si elle y est, l'accès est quasi instantané.

% 4. Registre
\subsection{Registre}
\begin{definitionbox}{Registre}
Les cases mémoire les plus rapides au monde, situées au cœur même des circuits de calcul de l'unité arithmétique du CPU.
\end{definitionbox}

\begin{analogybox}{Le Registre}
C'est la **main du cuisinier**. Il ne peut y tenir qu'un seul objet à la fois (la poêle ou le couteau) pour l'utiliser immédiatement.
\end{analogybox}

\begin{examplebox}{Le Registre}
Les registres \texttt{RAX}, \texttt{RBX} en architecture x86\_64 qui contiennent un nombre entier pendant une addition.
\end{examplebox}
\textbf{Pourquoi il existe :} Le CPU ne sait effectuer une addition qu'entre deux valeurs chargées directement dans ses registres.

% 5. Bus Système / Bus Mémoire
\subsection{Bus Système (Bus)}
\begin{definitionbox}{Bus}
Un ensemble de câbles électroniques partagés qui relient le CPU, les caches et la mémoire principale pour transporter les commandes et les données.
\end{definitionbox}

\begin{analogybox}{Le Bus}
C'est l'**interphone général de la cuisine** avec un haut-parleur. Quand quelqu'un parle dans l'interphone, toutes les personnes présentes dans la pièce l'entendent en même temps.
\end{analogybox}

\begin{examplebox}{Le Bus}
Le bus d'interconnexion ou le bus de données reliant les cœurs de processeurs à la mémoire centrale.
\end{examplebox}
\textbf{Pourquoi il existe :} Pour éviter de tirer des fils individuels entre chaque composant, tous se branchent sur cette "autoroute d'information" commune.\\
\textbf{Comment il est utilisé :} Quand le CPU a besoin d'une donnée, il crie sur le bus : \textit{"Envoyez-moi le contenu de l'adresse 00 !"}.

% 6. Lecture et Écriture
\subsection{Lecture (Read) et Écriture (Write)}
\begin{definitionbox}{Lecture (Read) \& Écriture (Write)}
\begin{itemize}
    \item **Lecture (Read)** : Action de consulter/recopier une valeur sans la modifier.
    \item **Écriture (Write)** : Action de modifier/remplacer la valeur d'une donnée.
\end{itemize}
\end{definitionbox}

\begin{analogybox}{Lecture \& Écriture}
\textit{Lire} c'est regarder l'heure sur sa montre (la valeur ne change pas). \textit{Écrire} c'est régler la montre à une nouvelle heure.
\end{analogybox}

% 7. Bloc Mémoire / Adresse / Cache Line
\subsection{Adresse Mémoire, Bloc Mémoire \& Ligne de Cache (Cache Line)}
\begin{definitionbox}{Adresse, Bloc et Ligne de Cache}
\begin{itemize}
    \item **Adresse Mémoire** : Le numéro unique d'une case mémoire (ex: Adresse \texttt{0x00}).
    \item **Bloc Mémoire / Ligne de Cache** : Le groupe fixe d'octets (ex: 64 octets) transféré en un seul bloc entre la RAM et le cache.
\end{itemize}
\end{definitionbox}

\begin{analogybox}{Bloc Mémoire}
Au lieu d'acheter un œuf individuel à l'épicerie, vous achetez une **boîte de 6 œufs**. La boîte est le bloc mémoire.
\end{analogybox}

% 8. Multi-cœur & Cohérence de Cache
\subsection{Multi-cœur \& Cohérence de Cache (Cache Coherence)}
\begin{definitionbox}{Cohérence de Cache}
Le mécanisme matériel qui garantit que lorsque plusieurs processeurs ont une copie d'un même bloc mémoire dans leurs caches respectifs, la modification du bloc par l'un soit immédiatement visible par les autres, sans qu'aucun ne lise une valeur périmée.
\end{definitionbox}

\begin{analogybox}{La Cohérence de Cache}
Imaginez 3 personnes ($P_1, P_2, P_3$) ayant chacune une fiche bristol dans leur poche (leur cache) avec la recette de la tarte aux pommes (donnée). Si $P_1$ décide de modifier la température de cuisson sur sa fiche, il doit crier l'information pour que $P_2$ et $P_3$ mettent à jour ou rayent leur propre fiche. Sinon, $P_2$ fera brûler la tarte !
\end{analogybox}

% 9. Read Hit, Read Miss, Write Hit, Write Miss
\subsection{Les Événements de Cache : Hit et Miss}
\begin{definitionbox}{Hit vs Miss}
\begin{itemize}
    \item **Read Hit (Succès de lecture)** : Le CPU veut lire une donnée et elle **est déjà présente et valide** dans son cache local.
    \item **Read Miss (Défaut de lecture)** : Le CPU veut lire une donnée mais elle **n'est pas** dans son cache (ou est invalide). Il doit la demander sur le bus.
    \item **Write Hit (Succès d'écriture)** : Le CPU veut écrire dans une donnée et il la possède **déjà valide** dans son cache.
    \item **Write Miss (Défaut d'écriture)** : Le CPU veut écrire dans une donnée qu'il **ne possède pas** dans son cache.
\end{itemize}
\end{definitionbox}

% 10. Invalidation, Write-Back
\subsection{Invalidation \& Write-Back}
\begin{definitionbox}{Invalidation \& Write-Back}
\begin{itemize}
    \item **Invalidation** : Action de marquer la copie d'un autre cache comme "périmée / obsolète" (Invalide).
    \item **Write-Back (Réécriture en mémoire)** : Action de recopier la valeur modifiée du cache local vers la mémoire principale (RAM) pour la rendre à jour.
\end{itemize}
\end{definitionbox}

% 11. Les 3 États du Protocole MSI : Invalid, Shared, Modified
\subsection{Les 3 États du Protocole MSI}
\begin{definitionbox}{États MSI}
\begin{itemize}
    \item **Invalide (I)** : La donnée dans ce cache n'est pas à jour ou n'existe pas. Inutilisable.
    \item **Partagé (Shared - S)** : La donnée est valide, propre (identique à la RAM), et peut être partagée par plusieurs processeurs en **lecture seule**.
    \item **Modifié (Modified - M)** : La donnée est valide, a été modifiée par ce processeur (différente de la RAM), et ce processeur est le **seul propriétaire exclusif**.
\end{itemize}
\end{definitionbox}

\newpage
% ==========================================
\section{Étape 3 — Construire Progressivement le Système}
% ==========================================

Assemblons la maison étape par étape !

\subsection{Étape 3.1 — Les Fondations : Architecture Uniprocesseur}
Dans un ordinateur simple avec un seul CPU et une mémoire RAM :
\begin{verbatim}
[ CPU ] <---> [ Cache L1 ] <---> [ Bus ] <---> [ Mémoire RAM ]
\end{verbatim}
Le CPU lit et écrit dans son cache. Si la donnée manque, il va en RAM via le Bus. Il n'y a aucun risque de conflit car le CPU est seul.

\subsection{Étape 3.2 — Les Murs : Architecture Multiprocesseur}
Ajoutons maintenant 3 processeurs ($P_1, P_2, P_3$), chacun doté de son propre cache privé L1, tous branchés sur le même Bus partagé et la même RAM :

\begin{verbatim}
[ CPU 1 ] <---> [ Cache P1 ] --+
                               |
[ CPU 2 ] <---> [ Cache P2 ] --+<===> [ BUS SYSTÈME ] <===> [ MÉMOIRE RAM ]
                               |
[ CPU 3 ] <---> [ Cache P3 ] --+
\end{verbatim}

\subsection{Étape 3.3 — Le Problème de Cohérence (La Torsion)}
Imaginons la séquence suivante :
\begin{enumerate}
    \item La RAM contient à l'adresse \texttt{00} la valeur \textbf{100}.
    \item $P_1$ lit l'adresse \texttt{00}. Sa valeur \textbf{100} est mise dans le Cache $P_1$.
    \item $P_2$ lit l'adresse \texttt{00}. Sa valeur \textbf{100} est mise dans le Cache $P_2$.
    \item Maintenant, $P_1$ écrit \textbf{999} à l'adresse \texttt{00} dans son Cache $P_1$.
    \item **Catastrophe !** Le Cache $P_2$ contient toujours **100** ! Si $P_2$ lit sa valeur, il lira une donnée fausse/obsolète !
\end{enumerate}

\subsection{Étape 3.4 — Le Toit : La Solution Snoopy (L'Écoute du Bus)}
Pour résoudre ce problème, les contrôleurs de cache intègrent un circuit appelé **Snoopy** (qui signifie \textit{"espion"}). 
\begin{itemize}
    \item Chaque cache **écoute en permanence** tout ce qui se raconte sur le Bus partagé.
    \item Quand $P_1$ veut écrire sur le Tag \texttt{00}, il crie sur le Bus : \textit{"Attention, j'invalide le Tag 00 !"}.
    \item Le contrôleur Snoopy de $P_2$ entend ce message sur le bus, repère qu'il possède aussi le Tag \texttt{00}, et transforme immédiatement son état en **Invalide** !
\end{itemize}

C'est cette logique qui est formalisée dans le schéma MEF !

\newpage
% ==========================================
\section{Étape 4 — Lecture Approfondie de l'Image (Machine à États Finis)}
% ==========================================

Analysons maintenant **chaque élément** de l'image (le schéma du protocole Snoopy Cache MSI).

\begin{center}
\textit{(Le schéma montre 3 bulles circulaires représentant les états d'une ligne de cache et les flèches de transition déclenchées par les actions CPU ou les requêtes du Bus.)}
\end{center}

% ------------------------------------------
\subsection{Les 3 Bulles (États de la Ligne de Cache)}
\begin{enumerate}
    \item **Bulle 1 : Invalide** : La ligne de cache ne contient aucune donnée utilisable.
    \item **Bulle 2 : Partagé (Shared)** : La donnée est propre, valide et consultable par le CPU local. D'autres CPU peuvent aussi posséder ce bloc en Partagé.
    \item **Bulle 3 : Modifié** : La donnée a été écrite localement. Elle est sale (plus récente que la RAM) et ce cache possède l'exclusivité totale.
\end{enumerate}

% ------------------------------------------
\subsection{Explication Détaillée de TOUTES les Flèches du Schéma}

\begin{transitionbox}{Flèche 1 : De Invalide ➔ Partagé (Top Left ➔ Top Right)}
\begin{itemize}
    \item **Événement déclencheur (en rouge)** : \texttt{CPU Read Miss}
    \item **Action sur le Bus (en noir)** : \texttt{Placer Read Miss sur le Bus}
    \item **Ce qui se passe dans le CPU** : Le CPU veut lire un bloc qui est à l'état \textit{Invalide} dans son cache.
    \item **Ce qui se passe dans le Cache** : Il ne peut pas servir la valeur immédiatement (Défaut).
    \item **Ce qui se passe sur le Bus** : Le contrôleur émet une requête \texttt{Read Miss (BusRd)} sur le bus partagé. Si un autre cache possédait le bloc à l'état \textit{Modifié}, celui-ci réalise un \textbf{Write-Back} vers la RAM. Le bloc est chargé et son état passe de \textbf{Invalide ➔ Partagé}.
\end{itemize}
\end{transitionbox}

\begin{transitionbox}{Flèche 2 : De Invalide ➔ Modifié (Top Left ➔ Bottom)}
\begin{itemize}
    \item **Événement déclencheur** : \texttt{CPU Write Miss}
    \item **Action sur le Bus** : \texttt{Placer write Miss sur le Bus}
    \item **Ce qui se passe** : Le CPU veut écrire dans un bloc qu'il ne possède pas (ou invalide). Il émet une requête d'écriture exclusive \texttt{BusRdX (Write Miss)} sur le bus. Tous les autres caches possédant ce bloc invalident leur copie. Le CPU écrit la nouvelle valeur dans son cache et le bloc passe de \textbf{Invalide ➔ Modifié}.
\end{itemize}
\end{transitionbox}

\begin{transitionbox}{Flèche 3 : Boucle sur Partagé (Shared ➔ Shared, En Haut à Droite)}
\begin{itemize}
    \item **Événement déclencheur** : \texttt{CPU Read Hit}
    \item **Action sur le Bus** : \textit{Aucune (Pas de message bus)}
    \item **Ce qui se passe** : Le CPU veut lire un bloc déjà présent à l'état \textit{Partagé}. Il lit la valeur directement dans son cache L1 à vitesse maximale. L'état reste **Partagé**.
\end{itemize}
\end{transitionbox}

\begin{transitionbox}{Flèche 4 : Autre Boucle / Re-demandé sur Partagé}
\begin{itemize}
    \item **Événement déclencheur** : \texttt{CPU Read Miss} (lors d'un remplacement de bloc)
    \item **Action sur le Bus** : \texttt{Placer Read Miss sur le Bus}
    \item **Ce qui se passe** : Le cache recharge le bloc en émettant une demande sur le bus. L'état reste **Partagé**.
\end{itemize}
\end{transitionbox}

\begin{transitionbox}{Flèche 5 : De Partagé ➔ Modifié (Top Right ➔ Bottom)}
\begin{itemize}
    \item **Événement déclencheur** : \texttt{CPU Write Hit}
    \item **Action sur le Bus** : \texttt{Placer write Invalidate sur le Bus}
    \item **Ce qui se passe** : Le CPU veut écrire dans un bloc qu'il possède déjà mais qui est en lecture seule (\textit{Partagé}). Pour devenir le seul propriétaire, il émet une demande d'invalidation \texttt{BusUpgr} sur le Bus. Les autres processeurs invalident leur copie. La donnée est modifiée localement et l'état passe de \textbf{Partagé ➔ Modifié}.
\end{itemize}
\end{transitionbox}

\begin{transitionbox}{Flèche 6 : De Partagé ➔ Modifié (Autre cas Write Miss)}
\begin{itemize}
    \item **Événement déclencheur** : \texttt{CPU Write Miss}
    \item **Action sur le Bus** : \texttt{Placer write Miss sur le Bus}
    \item **Ce qui se passe** : Le CPU écrit dans un bloc non présent en cache local, déclenchant l'invalidation globale et le passage direct de ce cache à l'état **Modifié**.
\end{itemize}
\end{transitionbox}

\begin{transitionbox}{Flèche 7 : Boucle sur Modifié (En Bas à Gauche)}
\begin{itemize}
    \item **Événement déclencheur** : \texttt{CPU Read Hit} OU \texttt{CPU Write Hit}
    \item **Action sur le Bus** : \textit{Aucune}
    \item **Ce qui se passe** : Le CPU lit ou réécrit dans un bloc qu'il possède déjà en exclusivité (\textit{Modifié}). Tout se fait localement en cache L1 sans encombrer le bus. L'état reste **Modifié**.
\end{itemize}
\end{transitionbox}

\begin{transitionbox}{Flèche 8 : Boucle Interne sur Modifié (Write Miss sur remplacement)}
\begin{itemize}
    \item **Événement déclencheur** : \texttt{CPU Write Miss}
    \item **Action sur le Bus** : \texttt{Write Back le bloc, Placer write Miss sur le Bus}
    \item **Ce qui se passe** : Si le CPU doit évincer un bloc sale pour faire de la place à un autre bloc écrit, il force la recopie de l'ancien bloc sale vers la RAM (\textbf{Write Back}) puis place son écriture sur le bus. L'état du nouveau bloc devient **Modifié**.
\end{itemize}
\end{transitionbox}

\begin{transitionbox}{Flèche 9 : De Modifié ➔ Partagé (Bottom ➔ Top Right)}
\begin{itemize}
    \item **Événement déclencheur** : \texttt{CPU Read Miss} (provenant d'un AUTRE processeur sur le Bus)
    \item **Action sur le Bus** : \texttt{Write Back le bloc, Placer Read Miss sur le Bus}
    \item **Ce qui se passe** : Le CPU $P_1$ possède un bloc à l'état \textit{Modifié}. Soudain, $P_2$ émet un \texttt{Read Miss} sur ce même bloc sur le Bus. Le contrôleur Snoopy de $P_1$ intercepte le message, effectue un **Write-Back** (écrit la valeur à jour dans la RAM), et rétrograde son état de \textbf{Modifié ➔ Partagé}.
\end{itemize}
\end{transitionbox}

\newpage
% ==========================================
\section{Étape 5 — Simulation Scénario Complet Pas à Pas}
% ==========================================

Mettons en scène une histoire complète avec **3 Processeurs ($P_1, P_2, P_3$)** et le **Bloc mémoire \texttt{00}** (valeur initiale en RAM = 50).

\subsection{Scénario 1 : Lecture Initiale \& Partage}
\begin{enumerate}
    \item **Étape 1 ($P_1$ veut lire \texttt{00})** :
    \begin{itemize}
        \item État initial de \texttt{00} dans Cache $P_1$ : \textbf{Invalide}.
        \item Déclencheur : \texttt{CPU Read Miss}.
        \item Action Bus : Émission de \texttt{Read Miss} sur le Bus.
        \item Résultat : Chargement de 50 depuis la RAM. État $P_1$ devient **Partagé (S)**.
    \end{itemize}
    
    \item **Étape 2 ($P_2$ veut aussi lire \texttt{00})** :
    \begin{itemize}
        \item État initial de \texttt{00} dans Cache $P_2$ : \textbf{Invalide}.
        \item Action Bus : Émission de \texttt{Read Miss} sur le Bus.
        \item Snooping : $P_1$ voit passer le message et reste à l'état **Partagé**.
        \item Résultat : $P_2$ charge 50. État $P_2$ devient **Partagé (S)**.
    \end{itemize}
\end{enumerate}

\subsection{Scénario 2 : Invalidation par Écriture}
\begin{enumerate}
    \item **Étape 3 ($P_1$ veut écrire la valeur \textbf{99} sur \texttt{00})** :
    \begin{itemize}
        \item État actuel de $P_1$ : \textbf{Partagé (S)}.
        \item Déclencheur : \texttt{CPU Write Hit}.
        \item Action Bus : $P_1$ émet \texttt{Write Invalidate} sur le Bus.
        \item Snooping par $P_2$ : Le contrôleur $P_2$ entend l'invalidation sur le Bus et passe son bloc \texttt{00} de **Partagé ➔ Invalide (I)**.
        \item Résultat : $P_1$ écrit 99 dans son cache et passe de **Partagé ➔ Modifié (M)**.
    \end{itemize}
\end{enumerate}

\subsection{Scénario 3 : Conflit \& Write-Back}
\begin{enumerate}
    \item **Étape 4 ($P_3$ veut lire l'adresse \texttt{00})** :
    \begin{itemize}
        \item État actuel de $P_3$ : \textbf{Invalide (I)}.
        \item Déclencheur : $P_3$ émet \texttt{Read Miss} sur le Bus.
        \item Snooping par $P_1$ : $P_1$ snoop le message et constate qu'il possède \texttt{00} à l'état **Modifié** (valeur 99).
        \item Action Write-Back : $P_1$ met en pause le bus, recopie sa valeur 99 dans la RAM centrale, puis rétrograde son état de **Modifié ➔ Partagé (S)**.
        \item Résultat : $P_3$ lit la valeur 99 mise à jour et passe à l'état **Partagé (S)**.
    \end{itemize}
\end{enumerate}

\newpage
% ==========================================
\section{Étape 6 — Résumé \& Pièges Fréquents aux Examens}
% ==========================================

\subsection{Résumé en 4 Points Clés}
\begin{enumerate}
    \item **Invalide (I)** : Pas de donnée valide $\rightarrow$ Déclenche un \textit{Miss} sur le Bus.
    \item **Partagé (S)** : Donnée propre (conforme à la RAM) et partagée $\rightarrow$ Permet des \textit{Read Hit} locaux sans bus.
    \item **Modifié (M)** : Donnée sale et exclusive $\rightarrow$ Permet des \textit{Write Hit} locaux sans bus.
    \item **Snooping** : L'écoute du bus garantit qu'une écriture invalide les autres caches ou force un Write-Back.
\end{enumerate}

\subsection{Les Erreurs les plus Fréquentes des Étudiants}
\begin{alertbox}{Erreur 1 : Oublier le Snooping}
Penser que les transitions dépendent uniquement du processeur local. \textbf{Faux !} Une transition de \textit{Modifié ➔ Partagé} est déclenchée par un \texttt{Read Miss} émis par un \textbf{autre} processeur sur le bus.
\end{alertbox}

\begin{alertbox}{Erreur 2 : Confondre Read Hit et Read Miss sur Partagé}
Un \texttt{Read Hit} sur l'état Partagé ne génère **aucun** trafic bus. Seul le \texttt{Read Miss} génère un message bus.
\end{alertbox}

\begin{alertbox}{Erreur 3 : Oublier le Write-Back}
Penser que la mémoire RAM est mise à jour immédiatement lors d'une écriture dans un cache en mode Write-Back. La RAM n'est mise à jour **que** lors de l'éviction ou lors du Snooping d'une lecture distante !
\end{alertbox}

\newpage
% ==========================================
\section{Étape 7 — Quiz de Vérification (10 Questions \& Corrigé)}
% ==========================================

\subsection{Questions}
\begin{enumerate}[label=\textbf{Q\arabic*:}]
    \item Quel est l'état d'un bloc de cache qui contient la seule version valide et modifiée d'une donnée ?
    \item Que se passe-t-il si un CPU fait un \texttt{Read Hit} sur un bloc à l'état \textit{Partagé} ?
    \item Quel message est envoyé sur le bus lors d'un \texttt{CPU Write Hit} sur un bloc à l'état \textit{Partagé} ?
    \item Si le processeur $P_1$ écrit dans un bloc à l'état \textit{Partagé}, quel sera l'état de ce bloc chez les autres processeurs ?
    \item Pourquoi dit-on que le protocole s'appelle \textit{Snoopy} ?
    \item Que réalise le contrôleur de cache lorsqu'il possède un bloc \textit{Modifié} et qu'il voit passer un \texttt{Read Miss} sur ce bloc ?
    \item Est-ce que plusieurs caches peuvent posséder le même bloc mémoire à l'état \textit{Modifié} en même temps ?
    \item Quelle est la différence entre la mémoire RAM et la mémoire Cache ?
    \item Si un bloc est à l'état \textit{Invalide}, le CPU peut-il y lire une donnée directement ?
    \item Qu'est-ce que le \textit{Write-Back} ?
\end{enumerate}

\newpage
\subsection{Corrigé Détaillé}

\begin{enumerate}[label=\textbf{Q\arabic*:}]
    \item **Réponse : L'état Modifié (M)**.\\
    \textit{Explication :} L'état Modifié garantit à la fois l'exclusivité (aucun autre cache n'a la donnée) et le fait qu'elle soit plus récente que la RAM.

    \item **Réponse : La donnée est lue directement en cache local sans aucun message sur le bus**.\\
    \textit{Explication :} Un Read Hit sur un bloc Partagé est une opération 100\% locale et ultra-rapide.

    \item **Réponse : Write Invalidate (BusUpgr)**.\\
    \textit{Explication :} Le CPU doit prévenir les autres caches de détruire leur copie avant de modifier la sienne.

    \item **Réponse : L'état Invalide (I)**.\\
    \textit{Explication :} Les autres caches reçoivent l'invalidation et marquent leur bloc comme périmé.

    \item **Réponse : Car les contrôleurs de cache "espionnent" (snoop) en permanence les messages circulant sur le bus partagé**.\\
    \textit{Explication :} Le terme vient du verbe anglais \textit{to snoop} (fureter/espionner).

    \item **Réponse : Il effectue un Write-Back vers la RAM et passe à l'état Partagé**.\\
    \textit{Explication :} Il doit fournir la version à jour à la RAM et au processeur demandeur.

    \item **Réponse : Non, c'est strictly impossible**.\\
    \textit{Explication :} L'état Modifié implique l'exclusivité absolue pour éviter les incohérences d'écriture.

    \item **Réponse : La RAM est grande mais plus lente ; le Cache est très petit mais ultra-rapide et situé près du CPU**.\\
    \textit{Explication :} Le cache sert de mémoire tampon haute vitesse pour le CPU.

    \item **Réponse : Non**.\\
    \textit{Explication :} L'état Invalide signifie que la valeur est périmée ou inexistante. Cela déclenche un Read Miss.

    \item **Réponse : L'action de réécrire la valeur modifiée d'un cache vers la mémoire principale (RAM)**.\\
    \textit{Explication :} Cela remet la mémoire principale à jour avec la dernière valeur modifiée.
\end{enumerate}

\end{document}
